% !TEX root = ../thesis.tex

\chapter{Analysis}\label{ch:analysis}

\section{Traditional methods of interpreting LiDAR data in archaeology}
Since the emergence of Airborne Laser Scanning technology in the 1970s~\cite{MILIARESIS20071076}, methods for interpreting data obtained using this technology have undergone a long evolution.
Although laser scanning has existed for decades, its broader adoption in archaeology and its widespread use for creating high-quality digital elevation models (DEM) began much later, mainly in the late 2000s.~\cite{HUERTAS1988131,zhang_unveiling_2024} This delay was caused not only by the need to process large point clouds produced by ALS, but also by the limited sensor accuracy of earlier systems and insufficient computational resources for detecting and analysing small archaeological microforms.~\cite{MILIARESIS20071076,michenot_identification_2024}
As a result, early archaeological remote sensing relied primarily on indirect visual cues rather than on a direct representation of bare-earth microtopography.~\cite{HUERTAS1988131,liow_use_1990} In this respect, the first generations of remote sensing-based interpretation differed substantially from the LiDAR-based workflows used today.~\cite{trier_using_2019}

\subsection{Remote Sensing before the widespread use of LiDAR data}
Before LiDAR became a routine source of archaeological terrain data, experts mainly worked with data sources that were informative, but limited in their ability to capture subtle terrain deformations. The most important of these were:
\begin{itemize}
    \item \textbf{Optical aerial photographs:} Since the 1930s, archaeological prospection has relied on visual analysis of photographs captured from aircraft flying over the investigated area.~\cite{cowley_what_2016} Because of the limited spatial resolution and the dependence on weather, season, and illumination conditions, these images were usually not sufficient for identifying small archaeological microforms directly. Instead, experts searched for indirect indicators such as crop marks~\cite{cowley_what_2016}, soil marks, and shadows cast by structures under favourable lighting conditions.~\cite{HUERTAS1988131,liow_use_1990}
    
    \item \textbf{Satellite data:} Programs such as Landsat~\cite{wulder_fifty_2022,williams_landsat_nodate}, declassified CORONA~\cite{ur_corona_2003,altmaier_digital_2002}, and Sentinel-1~\cite{michenot_identification_2024} provided access to large volumes of Earth observation data from many regions. However, low- and medium-resolution imagery was usually more suitable for identifying large landscape patterns than for detecting small archaeological features. In addition, optical satellite imagery remained strongly limited by vegetation cover, which often obscured the bare ground and reduced the visibility of terrain-based structures.~\cite{thompson_detecting_2020,michenot_identification_2024}
    
    \item \textbf{Early geometric 2D models:} In the 1980s and 1990s, early computer vision approaches operated mainly on 2D image features~\cite{griffiths_improving_2019} such as contours, angles, edges, and pixel intensity. In some cases, these methods attempted to infer 3D structure from shadows or simple geometric assumptions. Their applicability to archaeology was limited, because they worked best for regular man-made forms and were considerably less reliable for eroded, partially preserved, or topographically subtle remains.~\cite{HUERTAS1988131,liow_use_1990}
\end{itemize}

Taken together, these approaches were useful for detecting larger sites or visually distinctive anomalies, but they were much less effective when the goal was to identify faint archaeological relief features, especially in forested environments.~\cite{michenot_identification_2024} This limitation became one of the main reasons for the transition toward LiDAR-based terrain modelling, discussed in the following section.

\subsection{Transition to LiDAR and DEM}
Unlike passive optical sensors, LiDAR is an active sensing technology that can partially penetrate vegetation canopy and provide information about the terrain surface beneath it.~\cite{bundzel_semantic_2020,kokalj_standardizing_2025} With the development of point cloud filtering algorithms, it became possible to separate reflections from vegetation and other above-ground objects from ground returns and, on this basis, to construct bare-earth elevation models~\cite{MILIARESIS20071076,kokalj_standardizing_2025}. In archaeology, this represented a major shift, because it enabled researchers to analyse terrain morphology more directly, including subtle archaeological microrelief that is often invisible in conventional optical imagery~\cite{michenot_identification_2024,thompson_detecting_2020}. However, the production of a high-quality terrain model still depended on sufficiently dense point clouds and reliable interpolation and filtering procedures, which became more feasible only with later technological progress~\cite{kokalj_standardizing_2025,thompson_detecting_2020}.

At the same time, the term \emph{DEM} (Digital Elevation Model) is often used as a general designation for a digital representation of elevation values on a regular grid over a selected mathematical or physical surface~\cite{kokalj_standardizing_2025,zaksek_sky-view_2011}. In practice, point clouds acquired by ALS are transformed into raster models by interpolation, and several related model types are distinguished depending on which surface is represented.

The main model types relevant in this context are:
\begin{itemize}
    \item \textbf{DTM} (Digital Terrain Model) represents the terrain surface with vegetation and most above-ground objects removed~\cite{altmaier_digital_2002,kokalj_standardizing_2025}. In archaeological applications, such a model is highly valuable because it reveals the shape of the ground surface more clearly. At the same time, the distinction is not always terminologically strict in the literature, and \emph{the terms DEM and DTM are often used interchangeably}~\cite{kokalj_standardizing_2025}.
    
    \item \textbf{DSM} (Digital Surface Model) represents the upper surface including vegetation, buildings, and other objects above the ground~\cite{kokalj_standardizing_2025,altmaier_digital_2002}. For archaeological prospection in densely vegetated areas, DSM is generally less suitable, because it does not isolate the microtopography of the terrain itself.
    
    \item \textbf{DFM} (Digital Feature Model) may be used to denote a model that preserves terrain together with selected features of interest~\cite{kokalj_standardizing_2025}. In archaeological interpretation, this idea is relevant because researchers are often interested precisely in terrain-related anthropogenic features such as mounds, terraces, embankments, or remnants of ancient structures. However, terminology in the literature is not fully uniform, and \emph{many authors still refer to the working raster simply as DEM or DTM}~\cite{kokalj_standardizing_2025}.
\end{itemize}

Despite the major shift brought by high-resolution ALS data, the transition to semi-automated or fully automated interpretation methods was relatively slow~\cite{trier_using_2019,cowley_what_2016,verschoof-van_der_vaart_combining_2020}. Although LiDAR-derived terrain models provided more detailed and vegetation-reduced representations of the landscape, archaeological interpretation for a long time continued to rely primarily on expert-driven visual inspection~\cite{guyot_objective_2021,kokalj_standardizing_2025}. In other words, the data changed significantly, but the dominant interpretation workflow remained largely traditional.

One of the main approaches was the visual inspection of raster visualizations derived from the terrain model. Besides the base elevation model itself, researchers employed various visualizations designed to enhance subtle microrelief and improve the visibility of target archaeological features. The most commonly used visualizations include:
\begin{itemize}
    \item \textbf{Hillshade:} a shaded-relief visualization simulating illumination from a chosen direction, which helps reveal the shape of slopes and local terrain forms.~\cite{kokalj_standardizing_2025,horn_hill_1981,hobbs_investigation_1999}
    
    \item \textbf{Slope:} the first derivative of the elevation model, expressing the maximum rate of elevation change between neighbouring cells.~\cite{kokalj_standardizing_2025,bundzel_semantic_2020}
    
    \item \textbf{Sky-view factor (SVF):} a parameter expressing the fraction of visible sky from a given location; in archaeological visualization, it is useful because it reduces dependence on a single illumination direction.~\cite{zaksek_sky-view_2011,kokalj_standardizing_2025}
    
    \item \textbf{Openness:} a measure based on the angular openness of the horizon within a selected radius. In practice, positive openness is frequently used to emphasize convex relief forms and improve the detectability of archaeological features.~\cite{yokoyama_visualizing_2002,doneus_openness_2013,kokalj_standardizing_2025}
\end{itemize}

Within this workflow, the standard method of interpretation remained manual redrawing during visual inspection. An expert examined relief visualizations and manually delineated visible structures, often creating vectorized schematic representations that reflected the assumed shape, extent, and relative height of the observed archaeological objects~\cite{somrak_learning_2020,zhang_unveiling_2024}. This approach significantly improved archaeological mapping compared with earlier optical remote sensing, but it still depended strongly on expert experience, was time-consuming, and remained difficult to scale to large areas~\cite{cowley_what_2016,somrak_learning_2020,trier_using_2019}.

\subsection{Limitations of traditional approaches}
The traditional approaches described above have several limitations that hinder their systematic use and make large-scale application difficult. Although expert interpretation remains indispensable in archaeological remote sensing, manual inspection and annotation of LiDAR-derived relief visualizations are highly time-consuming and labour-intensive~\cite{somrak_learning_2020,zhang_unveiling_2024}. As the size of the investigated area increases, this workflow becomes progressively harder to maintain in a consistent manner~\cite{zhang_unveiling_2024,kokalj_standardizing_2025}. For this reason, large-scale and systematic archaeological mapping was long considered difficult to achieve using purely traditional expert-driven procedures~\cite{trier_using_2019}.

Another major limitation is that interpretation is inseparably linked to \emph{the cognitive framework of the researcher}~\cite{cowley_what_2016}. As emphasized in the literature, archaeological observation is conditioned by prior experience, familiarity with the data, and understanding of how the data were generated~\cite{cowley_what_2016,somrak_learning_2020}. In practice, this means that both the sensing technology and the \emph{intellectual framework of the interpreter} influence which features are noticed, how they are categorized, and how confidently they are recorded~\cite{cowley_what_2016}.

This problem becomes even more visible in long-term mapping projects and national archaeological databases, where classifications are often produced over many years by different researchers working with different objectives and procedures~\cite{cowley_what_2016}. Such heterogeneity may lead to inconsistencies in labeling, reduced comparability between regions, and difficulties in obtaining robust and reproducible results~\cite{cowley_what_2016,kokalj_standardizing_2025}.

For these reasons, \emph{the human factor became one of the main motivations} for research into semi-automatic and automatic approaches. These methods do not eliminate the role of the expert, but they can reduce the degree of subjective variability~\cite{somrak_learning_2020}, support more standardized analysis, and improve the efficiency of processing increasingly large ALS data sets~\cite{bundzel_semantic_2020,verschoof-van_der_vaart_combining_2020}.

\section{Automation and semi-automated methods}
Given the limitations of purely manual interpretation, the transition toward semi-automatic approaches was a natural development~\cite{trier_using_2019,blaschke_object_2010}. With the wider availability of DEM-based archaeological prospection, researchers began to develop methods based on spectral~\cite{gonzales2002digital} and geomorphometric analysis~\cite{wilson_terrain_2000}, image processing rules, and early computer vision techniques~\cite{longuet1982theory}. In contrast to later deep learning approaches, these methods relied mainly on explicitly defined mathematical rules and on archaeological domain knowledge about the expected morphology of target objects~\cite{verschoof-van_der_vaart_combining_2020,blaschke_object_2010}.

Most of these methods operate either on the numerical properties of terrain cells and their local neighbourhood, or on raster representations derived from terrain models~\cite{kokalj_standardizing_2025,blaschke_object_2010}. For the purposes of this review, they can be divided into two broad groups. The first group consists of \textit{indirect algorithms}, which operate primarily on pixel attributes such as elevation, slope, curvature, or local contrast~\cite{MILIARESIS20071076}. The second group consists of \textit{direct algorithms}, which analyse raster visualizations more explicitly as images~\cite{blaschke_object_2010,guyot_objective_2021}.

\subsection{Indirect algorithms}
In cases where raster visualizations did not provide sufficient visibility because of low contrast, noise, or complex terrain structure, a number of methods were developed to analyse the numerical properties of individual pixels in their local neighbourhood~\cite{MILIARESIS20071076}. Typical examples include:

\begin{itemize}
    \item \textbf{Region growing methods:} These approaches start from seed cells and iteratively merge neighbouring cells according to similarity criteria such as elevation, slope, or local terrain continuity~\cite{adams_seeded_1994,MILIARESIS20071076,zhu_loess_2018}. Their advantage is conceptual simplicity and the ability to exploit local terrain structure. Their weakness is high sensitivity to threshold selection and to local noise~\cite{adams_seeded_1994,castilla_size-constrained_nodate}.
    
    \item \textbf{Segmentation based on geomorphometric features:} In these approaches, terrain forms are represented by descriptors such as average height, slope, curvature~\cite{evans_geomorphometry_2012}, or shape-related parameters, after which clustering or rule-based classification is applied~\cite{MILIARESIS20071076,zhu_loess_2018}. Such methods made it possible to distinguish between broader classes of terrain forms, including potential anthropogenic features. However, their performance depended strongly on the quality of the selected descriptors and on how well archaeological objects could be represented by them~\cite{MILIARESIS20071076}.
    
    \item \textbf{Contour regularization and shape refinement:} Because initial segmentations were often noisy or geometrically unstable, additional post-processing procedures were applied to regularize contours~\cite{liow_use_1990}. These included straight-line approximation~\cite{douglas_algorithms_1973} and other knowledge-driven corrections intended to produce more plausible outlines of archaeological features~\cite{liow_use_1990,HUERTAS1988131}. While such procedures improved the visual quality of results, they also introduced additional assumptions and parameters that required expert tuning.
\end{itemize}

Overall, indirect methods were useful when archaeological objects could be characterized by relatively stable local terrain properties~\cite{guyot_objective_2021}. At the same time, they remained limited in heterogeneous landscapes, where natural and anthropogenic forms often overlap in their geomorphometric characteristics~\cite{trier_using_2019,bundzel_semantic_2020}.

\subsection{Direct algorithms}
Some terrain representations derived from DEM made archaeological features more visually separable from the surrounding background and therefore enabled more direct use of image-based analysis~\cite{kokalj_standardizing_2025,thompson_detecting_2020}. Although such methods did not completely replace indirect approaches, they provided a more intuitive bridge between expert visual interpretation and computer-assisted processing. The most representative approaches include:

\begin{itemize}
    \item \textbf{Template matching~\cite{brunelli2008template,brunelli2009template}:} This method was widely used for searching for discrete archaeological structures with relatively predictable shapes~\cite{trier_using_2019,guyot_objective_2021}. The algorithm scanned raster images for patterns corresponding to predefined templates or geometric forms~\cite{trier_automatic_2012}. Its main advantage was simplicity and direct applicability to visually distinctive objects. However, it performed poorly when object morphology varied substantially or when the background was heterogeneous, often resulting in a large number of false negatives and unstable performance across sites~\cite{trier_using_2019,guyot_objective_2021}.
    
    \item \textbf{Object-based image analysis (GEOBIA~\cite{Hay2008,blaschke_object_2010}):} This approach represented an important transition toward more structured automated interpretation~\cite{blaschke_object_2010,Castilla2008,chen_geographic_2018,hay_special_nodate}. Instead of classifying individual pixels, it first grouped pixels into meaningful image objects and then classified these objects using spatial, textural, and contextual descriptors such as area, compactness, or elongation~\cite{blaschke_object_2010,MILIARESIS20071076}. Compared with purely pixel-based methods, GEOBIA could better reflect the spatial organization of archaeological forms. Nevertheless, it still depended strongly on segmentation settings and manually selected descriptors, which limited its transferability between landscapes~\cite{castilla_size-constrained_nodate,diakogiannis_resunet-_2020}.
\end{itemize}

Although these approaches increased the efficiency of expert work, they still required substantial manual intervention. In practice, the analyst often had to adjust numerous parameters, such as contrast thresholds, filter radius, segmentation settings, or template definitions~\cite{somrak_learning_2020,kokalj_standardizing_2025}. As a result, these methods were difficult to scale reliably to large areas and remained sensitive to noise, relief complexity, and local landscape variability~\cite{zhang_unveiling_2024,bundzel_semantic_2020}. This limitation is one of the main reasons why later research increasingly shifted toward data-driven learning approaches, especially CNN~\cite{lecun_convolutional_1998}-based semantic segmentation~\cite{bundzel_semantic_2020,fiorucci_deep_2022}.

\section{Utilization of CNN and Deep Learning}
The shift toward deep learning methods in remote sensing and archaeological prospection is relatively recent~\cite{fiorucci_deep_2022,bundzel_semantic_2020}. Although convolutional neural networks (CNNs)~\cite{lecun_convolutional_1998} have older theoretical roots, their practical impact became significant mainly in the 2010s, when advances in computational power, optimization methods, and the availability of annotated datasets enabled the successful training of deep architectures for complex image analysis tasks~\cite{cui_semantic_2020,wurm_semantic_2019}. In archaeology, the interest in deep learning was motivated not only by its success in computer vision, but also by the practical limitations of expert-driven interpretation of large LiDAR-derived datasets, which is labor-intensive, time-consuming, and difficult to scale to regional or interregional analysis~\cite{somrak_learning_2020,verschoof-van_der_vaart_combining_2020}.

CNN-based approaches are particularly relevant for LiDAR-derived terrain representations because they can learn hierarchical features directly from data~\cite{chen_geographic_2018,zhao_semantic_2022}. In the first layers, the network typically captures local patterns such as edges, gradients, and small texture changes, whereas deeper layers encode more abstract spatial structures and contextual relationships~\cite{azarkhordad_detecting_2025,wurm_semantic_2019}. This is important for archaeological applications, where the target objects may differ in scale, preservation state, morphology, and surrounding terrain context~\cite{guyot_objective_2021,wurm_semantic_2019}. In contrast to traditional handcrafted approaches, CNNs do not rely exclusively on manually designed descriptors or fixed thresholds~\cite{wurm_semantic_2019,verschoof-van_der_vaart_combining_2020}. However, their performance remains strongly dependent on the quality of the training data, the consistency of annotation, and the suitability of the selected model for the specific task~\cite{bundzel_semantic_2020,griffiths_improving_2019}.

For clarity, it is necessary to distinguish the main computer vision tasks solved by deep learning models in remote sensing and archaeology:
\begin{itemize}
    \item \textbf{Semantic segmentation} assigns a class label to each pixel of the input image~\cite{cui_semantic_2020,chen_geographic_2018}. In the context of this thesis, it is the most directly relevant task, because it allows the creation of a dense mask indicating the spatial extent of archaeological objects such as terraces, mounds, or other terrain-related remains.
    \item \textbf{Object detection} localizes objects and assigns them a category, most commonly by predicting bounding boxes~\cite{fiorucci_deep_2022,azarkhordad_detecting_2025}. This task is useful when the goal is to identify the approximate position of discrete structures rather than their exact shape.
    \item \textbf{Instance segmentation} combines detection and segmentation by assigning a separate pixel mask to each individual object instance~\cite{bundzel_semantic_2020,guyot_objective_2021}. This is important when multiple objects of the same class appear close to one another and must be distinguished individually.
\end{itemize}

\subsection{Commonly used CNN architecture families}
For the purposes of this review, four broad architecture families are especially relevant in recent archaeological and remote sensing literature~\cite{fiorucci_deep_2022,zhang_unveiling_2024}. This categorization is not exhaustive, because many practical models combine components from multiple families and are further modified through different backbones, attention modules, loss functions, data augmentation strategies, and post-processing steps.

\textbf{Encoder--decoder architectures for semantic segmentation.}  
This family includes U-Net~\cite{ronneberger_u-net_2015} and its many derivatives, such as U-Net++~\cite{zhou_unetplusplus_2018}, ResUNet-a~\cite{diakogiannis_resunet-_2020,he_deep_2016}, and other modified encoder--decoder~\cite{zhang_semantic_2023} based architectures. Their key characteristic is the combination of a contracting path, which extracts increasingly abstract semantic features, with an expanding path, which restores spatial resolution~\cite{ronneberger_u-net_2015,bundzel_semantic_2020}. Skip connections between corresponding layers help preserve fine spatial detail and improve localization accuracy~\cite{ronneberger_u-net_2015,drozdzal_importance_2016}. This makes encoder--decoder architectures especially suitable for segmenting subtle archaeological terrain features in LiDAR-derived raster data, where both local detail and wider context are important~\cite{bundzel_semantic_2020,guyot_objective_2021}. In archaeological LiDAR research, U-Net has already shown strong applicability, and in the study by \textcite{bundzel_semantic_2020} it outperformed Mask R-CNN~\cite{he_mask_2018} for semantic segmentation of ancient Maya structures.

\textbf{DeepLab family.}  
DeepLabV3~\cite{chen_rethinking_2017} and DeepLabV3+~\cite{chen_encoder-decoder_2018} also follow the semantic segmentation paradigm~\cite{zhou_dynamic_2023,quan_improved_2021}, but they place stronger emphasis on multi-scale contextual modeling through atrous convolutions and the Atrous Spatial Pyramid Pooling (ASPP)~\cite{chen_rethinking_2017} module~\cite{zhou_dynamic_2023,liu_afnet_2021}. Their main advantage is the ability to enlarge the receptive field without excessive loss of spatial resolution~\cite{zhao_semantic_2022,quan_improved_2021}, which is useful when the target objects vary considerably in size or when broader terrain context is needed~\cite{liu_afnet_2021,diakogiannis_resunet-_2020} to distinguish archaeological structures from natural background forms. Although DeepLab-based models~\cite{quan_improved_2021,} are more frequently discussed in general remote sensing than in archaeological LiDAR specifically~\cite{zhang_unveiling_2024}, they remain an important comparative family because they address one of the central segmentation problems: combining local detail with broader contextual information~\cite{liu_afnet_2021,zhou_dynamic_2023}.

\textbf{Region-based CNN architectures.}  
This family~\cite{Girshick_2014_CVPR} includes Faster R-CNN~\cite{Girshick_2015_ICCV,ren_faster_2017} and Mask R-CNN~\cite{he_mask_2018}. These models are especially relevant for object detection and instance segmentation~\cite{fiorucci_deep_2022,guyot_objective_2021,verschoof-van_der_vaart_combining_2020}. Faster R-CNN~\cite{ren_faster_2017} is a two-stage detector that first generates candidate object regions and then classifies them~\cite{verschoof-van_der_vaart_combining_2020,griffiths_improving_2019,fiorucci_deep_2022}, while Mask R-CNN~\cite{he_mask_2018} extends this framework by adding a segmentation branch for each detected instance~\cite{guyot_objective_2021,fylaktos_deep_2025,griffiths_improving_2019,bundzel_semantic_2020}. In archaeological prospection, these architectures are useful when the primary goal is to identify and separate individual objects such as barrows, mounds, or building remains~\cite{fiorucci_deep_2022,verschoof-van_der_vaart_combining_2020,guyot_objective_2021}. They have already been applied to archaeological LiDAR data, including multi-class detection workflows and segmentation experiments~\cite{verschoof-van_der_vaart_combining_2020,guyot_objective_2021,bundzel_semantic_2020}. However, compared with encoder--decoder models, they are usually more naturally suited to discrete object instances than to dense semantic masks of continuous archaeological terrain modifications~\cite{bundzel_semantic_2020,fiorucci_deep_2022}.

\textbf{YOLO family and recent one-stage models.}  
YOLO~\cite{redmon_you_2016}-based models~\cite{ultralytics_yolo} were originally developed for fast object detection~\cite{griffiths_improving_2019}, but recent versions also support segmentation tasks~\cite{zhang_unveiling_2024}. Their main strength lies in speed and computational efficiency, which is attractive when very large LiDAR datasets need to be processed~\cite{zhang_unveiling_2024,azarkhordad_detecting_2025}. In archaeological research, their use is newer and still less established than U-Net- or R-CNN-based approaches~\cite{fiorucci_deep_2022,zhang_unveiling_2024}, but recent work has demonstrated promising results for the segmentation of Maya platforms and annular structures from aerial LiDAR imagery~\cite{zhang_unveiling_2024}. This suggests that YOLO-based models~\cite{ultralytics_yolo,zhang_unveiling_2024,azarkhordad_detecting_2025} may become increasingly relevant for practical large-scale archaeological workflows, especially when rapid screening is important~\cite{zhang_unveiling_2024,azarkhordad_detecting_2025}.

Across all these architecture families, the final result depends not only on the network design itself, but also on several related factors: the raster representation used as input~\cite{somrak_learning_2020,kokalj_standardizing_2025}, the quality and consistency of labels~\cite{griffiths_improving_2019,bundzel_semantic_2020}, the degree of class imbalance~\cite{Cui_2019_CVPR,zhou_dynamic_2023,bischke_segmentation_2018}, the use of transfer learning~\cite{pan_survey_2010,wurm_semantic_2019,zhang_semantic_2023}, and the chosen evaluation procedure~\cite{fiorucci_deep_2022,guyot_objective_2021}. For this reason, it would be incorrect to speak about a universally best CNN architecture for archaeological LiDAR analysis~\cite{bundzel_semantic_2020,somrak_learning_2020}. Instead, the choice should be matched to the specific task~\cite{fiorucci_deep_2022}. In the context of this thesis, which is focused on semantic segmentation of archaeological objects from airborne LiDAR-derived terrain data, encoder--decoder architectures are the most directly relevant family~\cite{zhao_semantic_2022}, while R-CNN- and YOLO-based approaches are valuable mainly as comparative alternatives or for tasks where individual object instances are of primary interest~\cite{zhang_unveiling_2024,verschoof-van_der_vaart_combining_2020}.

\subsection{What the authors agree on}
Despite the diversity of archaeological data sources, study areas, and model architectures, several recurring conclusions appear consistently across the literature.

\textbf{Transfer learning is often beneficial.}~\cite{cui_semantic_2020,wurm_semantic_2019}
A common observation is that transfer learning is highly useful when the available labeled dataset is limited, which is a typical situation in archaeological research~\cite{guyot_objective_2021,somrak_learning_2020}. Pre-trained weights, often originating from large natural-image datasets, can accelerate convergence, stabilize early training, and improve feature extraction in the low-data regime~\cite{cui_semantic_2020,zhang_semantic_2023}. At the same time, the literature does not suggest that transfer learning is a universal requirement in every setting. Rather, it is widely regarded as a practical strategy that often improves training efficiency and reduces the amount of archaeological training data needed~\cite{bundzel_semantic_2020,verschoof-van_der_vaart_combining_2020}.

\textbf{Data augmentation is an important part of the workflow.}~\cite{zhang_unveiling_2024,guyot_objective_2021}
Because archaeological datasets are usually relatively small and expensive to annotate, augmentation is widely used to improve model generalization and reduce overfitting~\cite{guyot_objective_2021,diakogiannis_resunet-_2020}. The literature most commonly highlights geometric transformations such as rotation, mirroring, scaling, and cropping, especially in cases where the target structures may appear in different orientations and local contexts~\cite{zhang_unveiling_2024,trier_using_2019}. In this sense, augmentation is not merely a technical convenience, but an important mechanism for improving robustness under limited-data conditions.

\textbf{Class imbalance is a central problem.}~\cite{zhang_unveiling_2024,zhou_dynamic_2023}
One of the most persistent difficulties in archaeological deep learning is the strong imbalance between target classes and background. In semantic segmentation, archaeological objects often occupy only a small fraction of all pixels~\cite{zhang_unveiling_2024,griffiths_improving_2019}, which can make optimization unstable and render some standard evaluation metrics misleading. As a result, many studies emphasize the need for imbalance-aware loss functions, careful sampling strategies, or alternative evaluation procedures that better reflect the performance on rare but important archaeological classes~\cite{azarkhordad_detecting_2025,fiorucci_deep_2022}.

\textbf{The quality of labels strongly affects model performance.}~\cite{griffiths_improving_2019,bundzel_semantic_2020}
Another point of broad agreement is that model quality is tightly connected to the quality of the reference annotations. Because archaeological labels are often produced through expert interpretation rather than direct physical boundaries observable in the data, they may contain uncertainty, simplification, or bias~\cite{cowley_what_2016,bundzel_semantic_2020}. These imperfections are then inherited by the trained model~\cite{rahaman_effects_2022,burgert_effects_2022}. For this reason, the preparation of training labels is increasingly treated as a methodological problem in its own right rather than as a trivial preprocessing step~\cite{Quintus_Day_Smith_2017,bundzel_semantic_2020}.

\subsection{Topics of discussion and disagreement}
Although there is broad agreement on the importance of deep learning for archaeological prospection, there is still no standardized workflow covering input representation, model design, data preparation, and evaluation~\cite{kokalj_standardizing_2025,fiorucci_deep_2022}. This leads to methodological variation and makes direct comparison between studies difficult.

One major topic of discussion concerns the \textbf{representation of LiDAR-derived data}. Archaeological deep learning models are rarely trained on raw elevation values alone. Instead, authors experiment with different visualizations or combinations of terrain derivatives, such as slope, sky-view factor~\cite{zaksek_sky-view_2011}, openness~\cite{yokoyama_visualizing_2002,doneus_openness_2013}, topographic position measures~\cite{weiss2001topographic}, or blended multichannel inputs~\cite{kokalj_standardizing_2025,zhang_unveiling_2024,bundzel_semantic_2020}. The underlying assumption is that an appropriate representation can improve the visibility of subtle archaeological features and facilitate learning. However, recent literature also points out that visualization practices remain insufficiently standardized, which reduces reproducibility and comparability across studies~\cite{kokalj_standardizing_2025}.

A second topic concerns the \textbf{choice of architecture}. While encoder--decoder networks are often preferred for semantic segmentation, region-based models and more recent one-stage detectors or segmenters may be advantageous in tasks focused on discrete object instances~\cite{fiorucci_deep_2022,bundzel_semantic_2020}. The literature therefore does not support a universally best architecture. Instead, the preferred model depends on the formulation of the task, the morphology of the target objects, the structure of the labels, and the scale of the study area~\cite{fiorucci_deep_2022,trier_using_2019}.

A third point of discussion is the \textbf{sampling strategy for training}. Because archaeological datasets are typically dominated by background, different studies adopt different approaches to the selection of positive and negative samples~\cite{verschoof-van_der_vaart_combining_2020,trier_using_2019}. Some workflows emphasize tiles containing objects of interest, whereas others deliberately preserve or reintroduce a limited number of negative examples in order to improve discrimination and reduce false positives~\cite{verschoof-van_der_vaart_combining_2020}. This indicates that sampling is not a purely technical detail, but one of the factors that can substantially influence the final model behavior.

A fourth unresolved issue concerns the \textbf{evaluation of results}. In remote sensing and archaeology, conventional pixel-wise metrics such as accuracy, IoU, precision, recall, and F1-score remain common~\cite{fiorucci_deep_2022,guyot_objective_2021}. However, several authors argue that these measures do not always reflect the practical usefulness of a model for archaeological interpretation~\cite{fiorucci_deep_2022,bundzel_semantic_2020,guyot_objective_2021}. For example, a prediction may be slightly misaligned at pixel level but still correctly identify an archaeological object in a way that is useful to an expert~\cite{fiorucci_deep_2022}. This has led to the proposal of alternative measures that better reflect archaeological reasoning, including centroid-based~\cite{fiorucci_deep_2022}, pixel-based~\cite{fiorucci_deep_2022}, and topology-aware~\cite{mosinska_beyond_2018} evaluation approaches~\cite{fiorucci_deep_2022,zhang_unveiling_2024,MILIARESIS20071076,grenier_accuracy_nodate}.

\subsection{The bottleneck of deep learning methods}
Despite the promising results achieved by deep learning methods, several critical limitations remain and currently form the main bottleneck of their wider use in archaeological research.

\textbf{Lack and uncertainty of reference data.}
The creation of high-quality annotated data is extremely labor-intensive and requires expert interpretation~\cite{somrak_learning_2020,bundzel_semantic_2020}. In archaeology, the label often does not correspond to a perfectly observable physical boundary, but rather to an expert-informed reconstruction of the likely extent of a feature~\cite{guyot_objective_2021,bundzel_semantic_2020}. This means that the resulting training data may include uncertainty~\cite{frenay_classification_2014}, simplification, and bias~\cite{cowley_what_2016,bundzel_semantic_2020}. Consequently, the model does not only learn archaeological patterns, but also the imperfections of the annotation process~\cite{rahaman_effects_2022}.

\textbf{Limited standardization of input representations.}
Another bottleneck lies in the lack of consensus regarding which LiDAR-derived representations should be used as model input~\cite{kokalj_standardizing_2025}. Since different visualizations highlight different aspects of topography, the final performance can depend strongly on representation choice. Without a more standardized and reproducible approach to visualization, it remains difficult to compare methods across projects and regions~\cite{kokalj_standardizing_2025}.

\textbf{Interpretability and expert trust.}
Deep learning models are still often criticized for their limited interpretability~\cite{azarkhordad_detecting_2025,bundzel_semantic_2020,}. In archaeology, this is particularly important because automated outputs are rarely accepted without contextual interpretation and, in many cases, field verification~\cite{fylaktos_deep_2025,guyot_objective_2021}. The difficulty of understanding why a model produced a specific prediction reduces trust in automatic findings and limits their direct adoption as independent evidence.

\textbf{Computational demands.}
Training modern deep learning models may require substantial computational resources, especially when dealing with large LiDAR-derived rasters, high-resolution imagery, or more complex architectures~\cite{fylaktos_deep_2025,zhang_semantic_2023}. This limits accessibility for smaller research teams and can constrain experimentation with model variants, loss functions, or extensive hyperparameter tuning~\cite{fylaktos_deep_2025}.

\textbf{Inference and post-processing issues.}
When large study areas are processed tile by tile, boundary artefacts and inconsistencies between neighboring predictions often emerge~\cite{zhang_unveiling_2024}. For that reason, inference usually requires additional strategies such as overlapping tiles, sliding-window prediction~\cite{dietterich_ensemble_2000}, or post-processing procedures to merge outputs into a coherent final map~\cite{bundzel_semantic_2020,cui_semantic_2020}.

\textbf{Mismatch between technical and archaeological evaluation.}
A final limitation is the gap between standard machine learning evaluation and archaeological usefulness~\cite{fiorucci_deep_2022}. High pixel-level performance does not automatically imply that the predictions are meaningful for archaeological interpretation, inventory, or decision support~\cite{fiorucci_deep_2022,bundzel_semantic_2020}. This is why recent research increasingly calls for evaluation protocols that better encode the logic of archaeological assessment and facilitate comparison between different workflows.

\section{The Expert's Role in Today's Research}
Despite the growing use of deep learning in archaeological prospection, the role of the expert has not diminished~\cite{bundzel_semantic_2020,fylaktos_deep_2025}. Rather, it has changed in character. The introduction of semi-automatic and automatic detection methods shifts expert work away from exhaustive manual search toward supervision, validation, and higher-level interpretation~\cite{verschoof-van_der_vaart_combining_2020,bundzel_semantic_2020}. In this sense, artificial intelligence does not remove the archaeologist from the workflow, but repositions the expert as a key actor~\cite{zanzotto_viewpoint_2019} in the formulation of training data, the assessment of model outputs~\cite{settles_active_2009}, and the integration of detected features into archaeological knowledge~\cite{bundzel_semantic_2020,fiorucci_deep_2022}.

Current deep learning workflows remain strongly dependent on the quality of the input data used for training and evaluation~\cite{griffiths_improving_2019}. In archaeology, this makes expert involvement particularly important, because the reference labels are rarely simple technical annotations~\cite{cowley_what_2016}. They are usually the result of archaeological interpretation informed by experience with LiDAR-derived terrain models, domain knowledge, and, where available, field verification~\cite{chase_airborne_2011,fylaktos_deep_2025,bundzel_semantic_2020}. For this reason, the preparation of training datasets is not a purely technical step, but one of the most methodologically sensitive parts of the entire workflow.

The expert contribution is especially visible in several critical processes:

\begin{itemize}
    \item \textbf{Creation of reference annotations:} The preparation of ground-truth data remains largely an expert-driven task~\cite{somrak_learning_2020,bundzel_semantic_2020}. Model convergence and final performance depend strongly on the quality, consistency, and interpretability of these annotations~\cite{azarkhordad_detecting_2025,rahaman_effects_2022,bundzel_semantic_2020}. At the same time, this process is one of the most labor-intensive stages of archaeological deep learning, often requiring substantial manual effort~\cite{somrak_learning_2020,bundzel_semantic_2020}.
    
    \item \textbf{Field verification and correction:} Archaeological interpretations derived from LiDAR data are often refined through fieldwork~\cite{fiorucci_deep_2022,bundzel_semantic_2020,chase_airborne_2011}. Ground verification helps clarify uncertain boundaries, confirm or reject proposed objects, and reduce systematic interpretation errors in the training data~\cite{fylaktos_deep_2025,bundzel_semantic_2020}.
    
    \item \textbf{Evaluation of model outputs:} The output of a deep learning model is not treated as final archaeological knowledge, but rather as a set of candidate detections or hypotheses~\cite{bundzel_semantic_2020,fylaktos_deep_2025}. These results must still be assessed by an expert in relation to terrain context, preservation state, and archaeological plausibility~\cite{zanzotto_viewpoint_2019}.
    
    \item \textbf{Historical and contextual interpretation:} Even when a model successfully identifies a feature, the archaeological meaning of that feature does not follow automatically from detection alone. Interpretation requires contextualization within settlement structure, land use, chronology, and wider socio-economic or cultural dynamics, which remains a fundamentally expert-led task~\cite{zanzotto_viewpoint_2019,azarkhordad_detecting_2025}.
\end{itemize}

Recent work increasingly shows that effective archaeological AI depends on close collaboration between archaeologists, geospatial data specialists, and machine learning researchers~\cite{olivier_implementing_2021,azarkhordad_detecting_2025,fiorucci_deep_2022}. This interdisciplinary cooperation is important not only for model development, but also for defining meaningful labels, selecting appropriate input representations, designing evaluation procedures, and interpreting the final outputs in ways that are relevant for archaeological research~\cite{olivier_implementing_2021}.

One of the main benefits of combining expert knowledge with trained models is enhanced perception. Deep learning systems can identify weak, repetitive, eroded, or spatially extensive patterns that may be difficult for a human observer to notice consistently across very large datasets~\cite{cowley_what_2016,azarkhordad_detecting_2025}. Conversely, the expert can critically assess these outputs, focus attention on ambiguous or unexpected cases, and determine whether they are archaeologically meaningful~\cite{fylaktos_deep_2025,bundzel_semantic_2020}. In this way, the model extends the expert's capacity rather than replacing expert judgment.

Another important consequence is methodological clarification~\cite{trier_using_2019,cowley_what_2016}. The use of machine learning often forces researchers to formulate more explicitly the criteria according to which archaeological objects are defined, labeled, and evaluated~\cite{trier_using_2019,bundzel_semantic_2020,azarkhordad_detecting_2025}. This can make interpretation workflows more transparent and reproducible than in purely intuitive visual inspection~\cite{kokalj_standardizing_2025,azarkhordad_detecting_2025,Castilla2008}.

Finally, model-assisted prioritization can improve the allocation of expert time and research resources~\cite{azarkhordad_detecting_2025,blaschke_object_2010,verschoof-van_der_vaart_combining_2020}. By highlighting areas or objects that are more likely to be relevant, deep learning methods can support more efficient planning of field verification and detailed archaeological analysis~\cite{fylaktos_deep_2025,zhang_unveiling_2024}.

Thus, in current archaeological practice, artificial intelligence should be understood primarily as a decision-support tool~\cite{azarkhordad_detecting_2025,bundzel_semantic_2020,fylaktos_deep_2025,zanzotto_viewpoint_2019,olivier_implementing_2021,shim_past_2002}. It can improve the localization and preliminary delineation of potential archaeological features, but validation, interpretation, and integration into broader historical understanding remain expert responsibilities~\cite{verschoof-van_der_vaart_combining_2020,guyot_objective_2021,azarkhordad_detecting_2025}.